
\section{Lower Bound}
\label{sec:lower}

\paragraph{Straight-line programs.}
Fix a field $\mathbb{F}$.
A (polynomial) straight-line program with input $x$ and parameters $p\in\mathbb{F}^m$ is a sequence of assignments that produces intermediate values $u_0,u_1,\dots$ and finally an output $P_p(x)$, where each assignment is obtained from earlier values (and from $x$, the parameters and fixed elements of $\mathbb{F}$) using additions and multiplications by fixed elements of $\mathbb{F}$, which are free, and a bounded number of \emph{nonscalar} multiplications, which are counted: a multiplication is nonscalar when neither factor is a fixed element of $\mathbb{F}$.
Thus multiplications by fixed scalars are free, but multiplications involving $x$ or the parameters are not.
Consequently each factor of a nonscalar multiplication, and the output, is a fixed linear combination of $x$ and the outputs of the earlier nonscalar multiplications plus an affine function of the parameters (the \emph{constant slot} of that factor); the parameters enter the program only through these constant slots.
An \emph{everywhere-defined rational left inverse} of a polynomial map $F\colon\mathbb{F}^m\to\mathbb{F}^m$ is a tuple of rational functions $g_i/h_i$ ($1\le i\le m$) such that no $h_i$ has a zero in $\mathbb{F}^m$ and $g_i(F(p))=p_i\,h_i(F(p))$ for every $p\in\mathbb{F}^m$.



\begin{lemma}[$6$ parameters]
\label{lem:six-params}
Assume $\mathbb{F}$ is infinite and $\mathrm{char}(\mathbb{F})\neq 2$.
Fix any six distinct evaluation points $x_0,\dots,x_5\in\mathbb{F}$, and fix a straight-line program as above that uses at most three nonscalar multiplications, has six scalar input parameters $p\in\mathbb{F}^6$, and outputs a polynomial $P_p(x)$.
Explicitly, the program is
\[
   u_1=(Ax+a)(Bx+b),\qquad u_2=\ell_2r_2,\qquad u_3=\ell_3r_3,\qquad
   P_p=s_3u_3+s_2u_2+s_1u_1+s_0x+b_1,
\]
where $A,B,s_i$ are fixed, $\ell_i,r_i$ are affine in $x$ and the earlier $u_j$ with fixed coefficients and constant slots $a_i,b_i$, and the seven constant slots $(a,b,a_2,b_2,a_3,b_3,b_1)$---both constants of the first multiplication included---are an arbitrary affine function of $p$.
Let
\[
   F:\mathbb{F}^6\to\mathbb{F}^6,\qquad
   F(p) = \bigl(P_p(x_0),\dots,P_p(x_5)\bigr).
\]
Then $F$ admits no everywhere-defined rational left inverse, and in particular no everywhere-defined rational inverse (equivalently: rational decoding with no exceptional inputs).
\end{lemma}

\paragraph{Connection to degree-$6$ decoding.}
A monic degree-$6$ polynomial is uniquely determined by its six remaining coefficients
\(
(c_5,\dots,c_0)\in\mathbb{F}^6
\),
and also by its values at the fixed six distinct points $x_0,\dots,x_5$ (since the corresponding Vandermonde matrix is invertible).
In the model investigated in this paper, ``rational decodability with no exceptional inputs'' means that the six degrees of freedom of the polynomial can be put in bijection (via everywhere-defined rational maps) with the six program parameters $p$.
Therefore we may reparameterize and assume $p=(c_5,\dots,c_0)$.
But then $F(p)$ is just evaluation at $x_0,\dots,x_5$, whose inverse is polynomial interpolation (the inverse Vandermonde transform), hence an everywhere-defined rational inverse.
The lemma rules this out.

\paragraph{Warm-up (Motzkin, degree $4$ with $2$ multiplications).}
The analogous statement is false for two multiplications.
Motzkin's classic program
\begin{align}
   u_1 &= x (x + a_0), \\
   u_2 &= (u_1 + x + a_1)(u_1 + a_2), \\
   P &= u_2 + a_3
\end{align}
uses exactly two nonscalar multiplications and produces a monic degree-$4$
polynomial.  Moreover its decoding is rational with no exceptional inputs when
$\mathrm{char}(\mathbb{F})\neq2$: comparing coefficients gives
\[
   c_3=2a_0+1,\quad
   c_2=a_0^2+a_0+a_1+a_2,\quad
   c_1=a_0(a_1+a_2)+a_2,\quad
   c_0=a_1a_2+a_3,
\]
so $a_0=(c_3-1)/2$; then $a_1+a_2$ is determined by $c_2$, whence
$a_2=c_1-a_0(a_1+a_2)$ and $a_1$, and finally $a_3=c_0-a_1a_2$.
Every step is a polynomial expression in $c_3,\dots,c_0$ with denominators
only powers of~$2$, so the Jacobian obstruction below cannot hold at degree~$4$.
Two multiplications therefore already evaluate all monic quartics; the
obstruction we prove is specific to degree~$6$ with three multiplications.

\paragraph{Motzkin's bound, and what is left to prove.}
Under \emph{any} preprocessing, a program with $m$ multiplications can produce at
most $2m$ independent coefficients: after the normalization of \Cref{appendix:lower}
(which trades the second constant of the first multiplication for an $x$-free
correction absorbed into the later constants), the first multiplication carries one
free slot, every later multiplication two, and the final addition one, so the
coefficient vector is a fixed polynomial function of $2m$ scalars, themselves
polynomial in the parameters, and its image cannot cover an $n$-dimensional family of monic polynomials
unless $2m\ge n$ (Motzkin~\cite{motzkin1955evaluation}; Theorem~M and exercise~30
in Knuth~\cite[\S4.6.4]{knuth1997seminumerical}, whose proof Knuth credits to
Pan~\cite{pan1966methods}; see also~\cite[Chapter~5]{burgisser1997algebraic}).
Hence over any infinite field, monic polynomials of degree $n$ need $\lceil n/2\rceil$
multiplications, whatever the preprocessing.
For odd $n$ this equals the $\lfloor n/2\rfloor+1$ of \Cref{thm:main}, so the
odd-degree schedules of this paper are optimal outright.
For even $n=2m$ the bound is $m$, and it is attained: by Motzkin and
Belaga~\cite{belaga1958some,pan1966methods} with algebraic preprocessing, and at
$n=6$ even with \emph{rational} preprocessing.  Pan's scheme~\cite{pan1961schemes}
(displayed as (16) in \cite[\S4.6.4]{knuth1997seminumerical})
\[
   z=(x+\alpha_0)x+\alpha_1,\qquad w=z+x+\alpha_2,\qquad
   P=\bigl((z-x+\alpha_3)\,w+\alpha_4\bigr)\,z+\alpha_5
\]
evaluates monic sextics with three multiplications, and its parameters are
rational functions of the coefficients---but with the denominator
$27c_3-18c_5c_4+5c_5^3$, so on that hypersurface the preprocessing is undefined.
(The bounds of Paterson and Stockmeyer~\cite{paterson1973evaluation} concern the
different model in which multiplications by the coefficients are free, where
$\Theta(\sqrt n)$ nonscalar multiplications are necessary and sufficient.)
The remaining question is therefore whether the exceptional set can be removed.
The decoders of this paper have no exceptional inputs (they divide only by
integers), and we conjecture that this costs exactly one multiplication at even
degrees: \emph{for $n\ge3$, evaluating monic polynomials of degree $2n$ with
everywhere-defined rational decoding requires at least $n+1$ multiplications}.
The restriction $n\ge3$ is necessary: the quartic program above evaluates every
monic degree-$4$ polynomial with two multiplications and rational decoding with no
exceptional inputs, so the statement fails at $n=2$.
Here we prove the first open case: degree $6$ cannot be done with three
multiplications, so Pan's hypersurface cannot be avoided.

\begin{proof}[Proof sketch]
The full proof is in \Cref{appendix:lower}; it is a Jacobian obstruction.
Let $J_F(p)$ be the $6\times6$ Jacobian of $F$.  If $F$ had an
everywhere-defined rational left inverse with coordinates $g_i/h_i$, then
$g_i(F(p))=p_i\,h_i(F(p))$ identically; differentiating at a point $p_0$ with
$J_F(p_0)v=0$, $v\neq0$, leaves $v_i\,h_i(F(p_0))=0$ for all $i$, which is
impossible because no $h_i$ vanishes.  So it suffices to exhibit, for every
three-multiplication program, one parameter point with singular Jacobian.

Every such program can be brought to a normal form
$u_1=x(R_{1,0}x+a_1)$, $u_2=\ell_2r_2$, $u_3=\ell_3r_3$ with $\ell_i,r_i$
affine in $x$ and the earlier $u_j$, and
$P=s_3u_3+s_2u_2+s_1u_1+s_0x+b_1$, by the gauge identity
$(Ax+a)(Bx+b)=A\,x(Bx+b+\frac BAa)+ab$, whose $x$-free correction $ab$ is
absorbed into the later constant slots.  The six normal-form slots
$q=(a_1,a_2,b_2,a_3,b_3,b_1)$ are therefore a polynomial map $q=Q(p)$ of
degree two of the parameters, and $J_F(p)=J(Q(p))\,DQ(p)$ with $J$ the
Jacobian of the normal form.  Unless $Q$ is an explicit quadratic
automorphism of $\mathbb{F}^6$, $DQ$ is singular at an explicit point (a
kernel direction of the slot map, or the midpoint of two parameter points
with the same normal-form slots, using $\mathrm{char}(\mathbb{F})\neq2$);
when it is, a singular point of $J$ pulls back to one of $J_F$.  So it
suffices to make $J$ singular for every normal form.  The sensitivities
$\partial P/\partial(\text{slot})$ are products of multiplicands with
the adjoints $\bar u_i=\partial P/\partial u_i$.  The argument splits on the
determinant $D$ of the coefficients with which $u_2,u_1$ enter the last gate.
If $D\neq0$, the slots can be chosen so that $\ell_3$ and $r_3$ take equal
values at two evaluation points while $\bar u_1,\bar u_2$ vanish there, so
two rows of $J$ coincide.  If $D=0$, a second-difference functional over
three evaluation points that annihilates constants and $x$ annihilates all
six sensitivities, so three rows of $J$ are dependent.  The normal-form
statement---both cases of this split, for every circuit in the normal form and
every affine reparameterization of its six slots---is machine-checked in Lean
(\texttt{FastPoly/LowerBound/}); the reduction to the normal form is not yet
formalized (see \Cref{appendix:lower}).
\end{proof}
\subsection{Beyond degree six}

We conjecture that the pattern continues: evaluating a generic monic
polynomial of degree $2n$ with everywhere-defined rational decoding requires at least $n+1$
multiplications.  The proof above suggests an approach: in the normal form,
each additional multiplication contributes two constant slots, hence two
Jacobian columns (and two further evaluation points, hence two rows), and the moment
functionals $\langle\,\cdot\,\rangle$ annihilating the partial derivatives
satisfy one quadratic (conic) consistency condition per level.  Beyond three
multiplications these conditions are no longer linear in the unknown
functional, and making the count rigorous is an open problem
(\Cref{sec:open-problems}).
